\chapter{Le jeu et ses exigences}
\label{chap-exigences}

\section{Règles}

Le joueur commande la flotte de la \emph{Sup Galilée Marine Nationale} contre une flotte de
pirates dirigée par l'ordinateur. Chaque camp dispose d'une grille de $10 \times 10$ cases et
d'une flotte de dix bateaux, soit vingt cases occupées (\cref{tab-flotte}).

\begin{table}[ht]
    \centering
    \begin{tabular}{@{}lccc@{}}
        \toprule
        Type de bateau  & Longueur (cases) & Nombre      & Cases occupées \\
        \midrule
        Quatre ponts    & 4                & 1           & 4              \\
        Trois ponts     & 3                & 2           & 6              \\
        Deux ponts      & 2                & 3           & 6              \\
        Un pont         & 1                & 4           & 4              \\
        \midrule
        \textbf{Flotte} &                  & \textbf{10} & \textbf{20}    \\
        \bottomrule
    \end{tabular}
    \caption{Composition d'une flotte.}
    \label{tab-flotte}
\end{table}

\begin{description}[style=nextline, leftmargin=1.2em]
    \item[Placement] Les bateaux sont droits, horizontaux ou verticaux. Deux bateaux ne se
          touchent jamais, \emph{ni par un côté, ni par un coin} : chaque bateau est entouré d'une
          ceinture d'eau (\cref{fig-non-contact}). Les flottes sont placées au hasard ; le joueur
          peut faire retirer la sienne tant qu'aucun coup n'a été tiré.
    \item[Tour de jeu] Le joueur commence. Tirer consiste à désigner une case de la grille
          adverse. Si un bateau s'y trouve, il est \emph{touché} et le tireur rejoue ; sinon le
          tir est \emph{manqué} et la main passe.
    \item[Bateau coulé] Un bateau dont toutes les cases sont touchées est \emph{coulé}.
    \item[Fin de partie] Le premier camp dont la flotte entière est coulée a perdu.
\end{description}

\begin{figure}[ht]
    \centering
    \begin{tikzpicture}[scale=0.62]
        % ------------------------------------------------ (a) placement licite
        \begin{scope}
            \fill[ecume] (1,2) rectangle (6,5);
            \plateau{7}
            \fill[marine, rounded corners=1.5pt] (2.12,3.12) rectangle (4.88,3.88);
            \fill[ocean, rounded corners=1.5pt] (6.12,0.12) rectangle (6.88,1.88);
            \fill[ocean, rounded corners=1.5pt] (0.12,6.12) rectangle (1.88,6.88);
            \node[note] at (3.5,-0.8) {(a) ceinture d'eau d'un bateau\\et voisins licites};
        \end{scope}
        % ---------------------------------------------- (b) placements interdits
        \begin{scope}[xshift=9.5cm]
            \plateau{7}
            \fill[marine, rounded corners=1.5pt] (2.12,3.12) rectangle (4.88,3.88);
            \fill[corail, rounded corners=1.5pt] (5.12,4.12) rectangle (5.88,5.88);
            \fill[corail, rounded corners=1.5pt] (1.12,2.12) rectangle (3.88,2.88);
            \draw[corail, line width=1.2pt] (4.75,3.75) -- (5.25,4.25) (4.75,4.25) -- (5.25,3.75);
            \draw[corail, line width=1.2pt] (2.75,2.75) -- (3.25,3.25) (2.75,3.25) -- (3.25,2.75);
            \node[note] at (3.5,-0.8) {(b) contact par un coin\\et contact par un côté};
        \end{scope}
    \end{tikzpicture}
    \caption{La règle de non-contact. En (a), la zone claire est interdite aux autres
        bateaux ; en (b), deux placements refusés.}
    \label{fig-non-contact}
\end{figure}

La règle de non-contact a une conséquence qui structure tout le projet : elle permet de
\emph{déduire}. Un coup au but prouve que ses quatre voisines diagonales sont de l'eau ; un
bateau coulé prouve que toute sa ceinture l'est aussi. Le jeu signale ces cases par un
drapeau, et la force de l'adversaire artificiel dépend d'abord de l'usage qu'il fait de ces
déductions (\cref{chap-adversaire}).

\section{Exigences fonctionnelles}

\begin{table}[ht]
    \centering
    \begin{tabularx}{\textwidth}{@{}lX@{}}
        \toprule
        \textbf{F1} & Jouer une partie complète contre l'ordinateur, selon les règles ci-dessus.   \\
        \textbf{F2} & Choisir parmi trois niveaux de difficulté : facile, moyen, difficile.        \\
        \textbf{F3} & Placer les flottes au hasard ; retirer celle du joueur avant le premier tir. \\
        \textbf{F4} & Afficher l'état des deux grilles : bateaux du joueur, tirs manqués, coups
        au but, bateaux coulés, cases exclues par déduction.                                       \\
        \textbf{F5} & Afficher le score de chaque camp et le tour en cours.                        \\
        \textbf{F6} & Consulter les règles et la présentation du jeu à tout moment.                \\
        \textbf{F7} & En fin de partie, proposer de rejouer, de changer de niveau ou de quitter.   \\
        \textbf{F8} & Jouer au clavier comme à la souris ; effets sonores facultatifs.             \\
        \bottomrule
    \end{tabularx}
    \caption{Exigences fonctionnelles.}
    \label{tab-fonctionnelles}
\end{table}

\section{Exigences de conception}

Ce sont elles qui distinguent la version~2 de la première. Chacune est vérifiable, et le
\cref{chap-qualite} montre par quel moyen elle l'est effectivement.

\begin{table}[ht]
    \centering
    \begin{tabularx}{\textwidth}{@{}lXl@{}}
        \toprule
                    & \textbf{Exigence}                                                        & \textbf{Vérifiée par} \\
        \midrule
        \textbf{C1} & Le modèle ne dépend ni de GTK ni d'aucune autre bibliothèque.            &
        compilation sans GTK                                                                                           \\
        \textbf{C2} & Une même graine rejoue exactement la même partie, sur toute plate-forme. &
        fichier de référence                                                                                           \\
        \textbf{C3} & L'adversaire ne peut pas tricher : il ignore la position des bateaux.    &
        construction (types)                                                                                           \\
        \textbf{C4} & Une opération refusée laisse l'état du jeu intact.                       & tests unitaires       \\
        \textbf{C5} & Aucun avertissement du compilateur ; aucune erreur mémoire.              &
        \code{-Werror}, ASan, UBSan                                                                                    \\
        \textbf{C6} & Le programme se lance depuis n'importe quel dossier.                     &
        ressources incorporées                                                                                         \\
        \textbf{C7} & L'interface ne se fige jamais pendant le tour de l'ordinateur.           &
        un tir par tic d'horloge                                                                                       \\
        \bottomrule
    \end{tabularx}
    \caption{Exigences de conception et moyens de vérification.}
    \label{tab-conception}
\end{table}
